% Chapter 15: quant-lib — The Unified Facade
% Companion code: crates/quant-lib

\chapter{quant-lib: The Unified Facade}
\label{ch:quant-lib}

\begin{companioncode}{quant-lib}
Code: \texttt{crates/quant-lib/}\\
Dependencies: all source crates (re-exported)\\
Run: \texttt{cargo run -p quant-lib --example 15\_full\_pipeline}\\
Run: \texttt{cargo run -p quant-lib --example 16\_trait\_compose}
\end{companioncode}

\section{Learning Objectives}

This chapter consolidates the ecosystem built in Chapters 6--14 into a
single facade crate. After completing it you can:

\begin{itemize}
    \item Understand the facade pattern as the final consolidation step
    \item Use the module structure of \crate{quant-lib} and its re-exports of the underlying crates
    \item Use the prelude for ergonomic one-line imports
    \item Compose pluggable components via the \texttt{BacktestBuilder} fluent API
    \item Run the end-to-end pipeline that ties five crates together
\end{itemize}

\section{Introduction: From Crates to a Library}

Chapters 1 through 14 built nine specialised crates: \crate{quant-core},
\crate{quant-timeseries}, \crate{quant-vol}, \crate{quant-stochastic},
\crate{quant-options}, \crate{quant-portfolio}, \crate{quant-factors},
\crate{quant-microstructure}, and \crate{quant-backtest}. Each crate
is independently useful, but a downstream user who wants the whole
ecosystem would rather depend on one crate than nine. \crate{quant-lib}
is that one crate.

\begin{keyinsight}
\crate{quant-lib} is purely a re-export facade. It adds \emph{zero}
new logic. Every line in its source is \texttt{pub mod} or
\texttt{pub use}. The math, types, and traits all live in the
underlying crates. The facade exists to give downstream projects a
single dependency and a unified namespace.
\end{keyinsight}

\section{The Facade Pattern}

A \emph{facade} is a structural pattern: a single object (here, a
crate) that provides a simplified, unified interface to a larger body
of code. The facade does not implement new behaviour; it delegates.

In Rust, the facade crate is structurally tiny:

\begin{lstlisting}[language=Rust]
// crates/quant-lib/src/lib.rs (excerpt)
pub mod prelude;

#[cfg(feature = "core")]
pub mod core {
    pub use quant_core::*;
}

#[cfg(feature = "options")]
pub mod options {
    pub use quant_options::*;
}

#[cfg(feature = "backtest")]
pub mod backtest {
    pub use quant_backtest::*;
}

// ... other modules ...
\end{lstlisting}

Each module is a thin re-export. Downstream code can write either
\texttt{use quant\_core::mean} or \texttt{use quant\_lib::core::mean}---
both resolve to the same function. The facade adds a layer of
indirection that costs nothing at runtime (everything is monomorphised
away) but gives the user a single dependency to track.

\section{Module Map}

\crate{quant-lib} organises the ecosystem into ten modules, one per
underlying crate plus a dedicated \module{traits} module for the
pluggable component trait catalogue.

\begin{center}
\begin{tabular}{lll}
\toprule
\textbf{Module} & \textbf{Source crate} & \textbf{Contents} \\
\midrule
\func{core}          & \crate{quant-core}          & moments, rolling windows, RNG, GBM \\
\func{traits}        & \crate{quant-core}          & Pluggable component traits \\
\func{risk}          & \crate{quant-core}          & PSR, DSR, Calmar, Omega, Ulcer \\
\func{timeseries}    & \crate{quant-timeseries}    & OLS, ACF, ADF, frac-diff, CUSUM \\
\func{vol}           & \crate{quant-vol}           & EWMA, ARCH, GARCH \\
\func{stochastic}    & \crate{quant-stochastic}    & Brownian, GBM, Monte Carlo \\
\func{options}       & \crate{quant-options}       & Black-Scholes, Greeks, IV \\
\func{portfolio}     & \crate{quant-portfolio}     & Markowitz, tangency, CAPM, VaR \\
\func{factors}       & \crate{quant-factors}       & PCA, Fama-French, attribution \\
\func{microstructure}& \crate{quant-microstructure}& LOB, OFI, impact, VWAP \\
\func{backtest}      & \crate{quant-backtest}      & triple-barrier, purged CV, Kelly \\
\bottomrule
\end{tabular}
\end{center}

\section{The Prelude}

The prelude pulls in the most commonly used items from every module
so that user code can start with a single import:

\begin{lstlisting}[language=Rust]
use quant_lib::prelude::*;
\end{lstlisting}

The prelude is intentionally curated. It includes:
\begin{itemize}
    \item Core utilities: \func{mean}, \func{std\_dev}, \func{variance},
          \func{RollingWindow}, \func{XorShift64}, \func{PriceSeries}.
    \item Pluggable component traits: \func{CrossValidator}, \func{Labeler},
          \func{BetSizer}, \func{StochasticProcess}, \func{OptionPricer},
          \func{Greeks}, \func{FactorModel}, \func{ImpactModel},
          \func{OrderBookOps}, \func{SampleWeighter},
          \func{StructuralBreakDetector}.
    \item Risk metrics: \func{probabilistic\_sharpe\_ratio},
          \func{deflated\_sharpe\_ratio}, \func{calmar\_ratio},
          \func{omega\_ratio}, \func{information\_ratio},
          \func{ulcer\_index}.
    \item Time series: \func{ols}, \func{adf\_test}, \func{fracDiff},
          \func{OlsFit}.
    \item Volatility: \func{ewma\_vol}, \func{ArchModel},
          \func{GarchModel}.
    \item Stochastic: \func{gbm}, \func{brownian\_motion},
          \func{bs\_call}, \func{bs\_put}, \func{mc\_call},
          \func{mc\_put}.
    \item Options: \func{delta}, \func{gamma}, \func{vega},
          \func{theta}, \func{implied\_vol}.
    \item Portfolio: \func{Portfolio}, \func{efficient\_frontier\_point},
          \func{min\_variance\_portfolio}, \func{tangency\_portfolio},
          \func{sharpe\_ratio}.
    \item Factors: \func{pca}, \func{ff3\_regression},
          \func{risk\_attribution}, \func{PcaResult},
          \func{FF3Exposure}.
    \item Microstructure: \func{OrderBook}, \func{Side}, \func{vwap},
          \func{sqrt\_impact}, \func{execution\_cost}.
    \item Backtest: \func{GenericBacktest}, \func{BacktestBuilder},
          \func{TripleBarrierLabel}, \func{TripleBarrierConfig},
          \func{FixedHorizonLabeler}, \func{KellyBetSizer},
          \func{WalkForward}, \func{WalkForwardConfig},
          \func{PurgedKFoldConfig}, \func{triple\_barrier\_label},
          \func{purged\_kfold\_splits}, \func{sample\_weights},
          \func{kelly\_fraction}, \func{walk\_forward\_efficiency},
          \func{afml\_backtest}.
\end{itemize}

Items not in the prelude are still reachable through the module path,
e.g., \texttt{quant\_lib::backtest::EqualBetSizer}. The prelude is a
convenience, not a gate.

\section{Feature Flags for Module Gating}

The \crate{quant-lib} \texttt{Cargo.toml} exposes each module behind a
feature flag so users pay only for what they compile:

\begin{lstlisting}
[features]
default = ["all"]
all = ["core", "timeseries", "vol", "stochastic",
        "options", "portfolio", "factors",
        "microstructure", "backtest"]
\end{lstlisting}

A user who needs only option pricing can disable defaults and pick
two modules:

\begin{lstlisting}
[dependencies]
quant-lib = { version = "0.1", default-features = false, features = ["core", "options"] }
\end{lstlisting}

The rest of the crate is not compiled. In a CI pipeline that builds
only the pricing surface, this shaves minutes off the build.

\section{Root-Level Convenience Re-exports}

Beyond the module re-exports, \crate{quant-lib} also exposes a small
set of root-level convenience functions so that the most common
operations do not require a module prefix:

\begin{lstlisting}[language=Rust]
// These three are equivalent:
quant_lib::core::mean(&data);
quant_lib::mean(&data);           // root re-export
quant_core::mean(&data);          // direct from source

// Same for stochastic, options, portfolio, backtest:
quant_lib::bs_call(s, k, r, sigma, t);
quant_lib::delta(s, k, r, sigma, t);
quant_lib::tangency_portfolio(&cov, rf);
quant_lib::kelly_fraction(p, b, a);
\end{lstlisting}

The root re-exports exist because these functions account for the
overwhelming majority of calls in examples and tests. Putting them at
the root saves typing without obscuring the module structure.

\section{Composable Backtest Architecture}

The headline feature of the trait architecture is the \func{GenericBacktest}
engine, generic over three pluggable components:

\begin{lstlisting}[language=Rust]
pub struct GenericBacktest<L, CV, BS>
where
    L: Labeler,
    CV: CrossValidator,
    BS: BetSizer,
{
    labeler: L,
    validator: CV,
    sizer: BS,
    entry_step: usize,
}
\end{lstlisting}

The bound \texttt{where L: Labeler, CV: CrossValidator, BS: BetSizer}
is the entire contract. Adding a new labeler means writing an
\texttt{impl Labeler for MyLabeler}---no engine changes, no new
features, no feature flags.

\subsection{The Builder API}

The \func{BacktestBuilder} gives a fluent API for composition:

\begin{lstlisting}[language=Rust]
let bt = BacktestBuilder::new()
    .labeler(FixedHorizonLabeler::new(5, 0.01))
    .cv(WalkForward::new(WalkForwardConfig::rolling(100, 25, 25)))
    .sizer(KellyBetSizer::new(0.5))   // half Kelly
    .entry_step(5)
    .build();

let results = bt.run(&prices, &returns)?;
\end{lstlisting}

Swapping one component is a one-line change. Example~16 in the
companion code runs the same strategy with four different bet sizers
(full Kelly, half Kelly, fixed 10\%, equal 100\%) and two different
validators (rolling vs.\ anchored walk-forward) by changing one
argument in the builder call.

\subsection{What the Builder Checks}

The builder rejects incomplete configurations at \emph{build time},
returning \func{BacktestError}:

\begin{itemize}
    \item Missing labeler: \texttt{BacktestError::MissingLabeler}.
    \item Missing cross-validator: \texttt{BacktestError::MissingCrossValidator}.
    \item Missing bet sizer: \texttt{BacktestError::MissingBetSizer}.
    \item Entry step out of range: \texttt{BacktestError::InvalidEntryStep}.
\end{itemize}

This means misconfigurations are caught before any data is loaded,
not in the middle of a 252-day backtest.

\section{The End-to-End Pipeline}

Example~15 in the companion code ties together five crates in a
single pipeline:

\begin{lstlisting}[language=Rust]
use quant_lib::prelude::*;
use quant_lib::backtest::{compute_position_size, kelly_from_returns};
use quant_lib::core::simple_returns;

// 1. Simulate a GBM price path (quant-stochastic).
let mut rng = XorShift64::new(123);
let prices = gbm(100.0, 0.08, 0.20, 1.0, 252, &mut rng);

// 2. Price an ATM call with the trait-based BlackScholes pricer.
let pricer = quant_options::BlackScholes::new(0.05, 0.20);
let call = pricer.price(100.0, 100.0, 1.0, OptionType::Call)?;

// 3. Analytical Greeks via the Greeks trait.
let d = pricer.delta(100.0, 100.0, 1.0, OptionType::Call)?;
let g = pricer.gamma(100.0, 100.0, 1.0)?;
let v = pricer.vega(100.0, 100.0, 1.0)?;
let th = pricer.theta(100.0, 100.0, 1.0, OptionType::Call)?;

// 4. Kelly fraction from synthetic trade returns.
let rets = simple_returns(&prices);
let f_kelly = kelly_from_returns(&rets);

// 5. Walk-forward cross-validation on the simulated path.
let wf = WalkForward::new(WalkForwardConfig::rolling(120, 30, 30));
let splits = wf.splits(0, prices.len());

// 6. Run the generic composable backtest.
let bt = GenericBacktest::new(
    FixedHorizonLabeler::new(5, 0.01),
    WalkForward::new(WalkForwardConfig::rolling(100, 25, 25)),
    KellyBetSizer::new(0.5),
    5,
);
let results = bt.run(&prices, &bt_returns)?;
\end{lstlisting}

\begin{companioncode}{quant-lib}
See \texttt{crates/quant-lib/examples/15\_full\_pipeline.rs} and
\texttt{crates/quant-lib/examples/16\_trait\_compose.rs}. Run with
\texttt{cargo run -p quant-lib --example 15\_full\_pipeline}.
\end{companioncode}

This is the ``hello world'' of the facade: every step uses a different
sub-crate, all reached through \crate{quant-lib}. The pipeline
demonstrates that the trait-based architecture composes cleanly---the
\func{BlackScholes} pricer implements \func{OptionPricer} and
\func{Greeks}; the \func{FixedHorizonLabeler} implements
\func{Labeler}; the \func{WalkForward} validator implements
\func{CrossValidator}; the \func{KellyBetSizer} implements
\func{BetSizer}. The \func{GenericBacktest} engine stitches them
together without knowing any of their concrete types.

\section{The 16-Example Tour}

\crate{quant-lib} ships 16 numbered examples that progress from
single-crate usage to multi-crate composition:

\begin{center}
\begin{tabular}{rlp{6.5cm}}
\toprule
\# & \textbf{Example} & \textbf{Crates exercised} \\
\midrule
01 & mean\_variance           & core \\
02 & returns\_basics          & core \\
03 & rolling\_window          & core \\
04 & random\_walk             & core, stochastic \\
05 & ols\_regression          & core, timeseries \\
06 & stationarity             & timeseries \\
07 & volatility               & vol \\
08 & options\_greeks          & options \\
09 & portfolio\_frontier      & portfolio \\
10 & pca\_factors             & factors \\
11 & microstructure\_lob      & microstructure \\
12 & triple\_barrier          & backtest \\
13 & purged\_cv               & backtest \\
14 & kelly\_sizing            & backtest \\
15 & full\_pipeline           & stochastic, options, backtest \\
16 & trait\_compose           & backtest (builder pattern) \\
\bottomrule
\end{tabular}
\end{center}

Examples 1--14 each exercise a single sub-crate. Example~15 is the
multi-crate pipeline. Example~16 is the composable backtest with four
bet sizers and two validators. The progression is deliberate: each
example introduces one new concept, so a reader working through them
in order never encounters a dependency they have not already seen.

\section{Non-Destructive by Design}

\crate{quant-lib} is a strict facade. The source crates continue to
work without \crate{quant-lib} in their dependency tree:

\begin{lstlisting}[language=bash]
# quant-backtest still builds without quant-lib:
cargo test -p quant-backtest
\end{lstlisting}

The 20-test TDD contract for \crate{quant-lib} (Appendix~B lists the
full suite) includes a test that verifies this property: the source
crates must pass independently. This guarantees that the facade
remains a convenience, not a coupling mechanism. A downstream project
that depends on \crate{quant-backtest} directly, without
\crate{quant-lib}, is not broken when \crate{quant-lib} adds a new
re-export.

\begin{keyinsight}
The facade is a one-way dependency: \crate{quant-lib} depends on the
source crates, never the reverse. If you find yourself writing
\texttt{use quant\_lib} inside a source crate, stop---you are
creating a cycle. The source crates must remain oblivious to the
facade's existence.
\end{keyinsight}

\section{Versioning Strategy}

When a source crate bumps its version (e.g.,
\crate{quant-options} from 0.1 to 0.2), \crate{quant-lib} bumps its
own version and updates the \texttt{Cargo.toml} dependency. Downstream
code that uses \crate{quant-lib} sees a single version bump, not
nine. This is the main practical benefit of the facade for consumers:
the dependency graph has one node, not nine.

\section{When to Use quant-lib vs.\ Direct Crates}

\begin{center}
\begin{tabular}{lp{8cm}}
\toprule
\textbf{Use} & \textbf{When} \\
\midrule
\crate{quant-lib}        & You want the whole ecosystem, one dependency, and unified imports. \\
\crate{quant-core}       & You need only moments, rolling windows, or the trait catalogue. \\
\crate{quant-options}    & You are building a pricing tool and do not need backtesting. \\
\crate{quant-backtest}   & You are building a research tool and do not need pricing. \\
individual crates        & You want minimum compile time and a small dependency tree. \\
\bottomrule
\end{tabular}
\end{center}

There is no wrong answer. The facade and the source crates are both
first-class supported surfaces. The choice is about ergonomics versus
minimality.

\section{Exercises}

\begin{enumerate}
    \item \textbf{Module discovery.} List the public items re-exported
          by \crate{quant-lib::portfolio}. How many of them are
          types, how many are functions, and how many are traits?
    \item \textbf{Prelude contents.} The prelude does not include
          \func{EqualBetSizer} or \func{FixedBetSizer}. Why? Add them
          to your own local prelude and rebuild---do any examples
          break?
    \item \textbf{Feature flags.} Build \crate{quant-lib} with
          \texttt{default-features = false, features = ["core",
          "options"]} and run example 08. Which modules are
          unavailable? What is the compile-time difference?
    \item \textbf{Builder swap.} Modify example 16 to use a
          \func{TripleBarrierLabeler} instead of
          \func{FixedHorizonLabeler}. What single line do you change?
          Does the rest of the pipeline still compile?
    \item \textbf{Pipeline extension.} Extend example 15 to compute
          the implied volatility from the simulated option price
          using \func{implied\_vol}. How many lines do you add? Do
          you need to change any existing line?
    \item \textbf{Root re-exports.} Identify three functions in the
          root re-export list that you would remove if you were
          designing the facade for a minimal API. Justify each
          choice.
    \item \textbf{Non-destructive check.} Run \texttt{cargo test -p
          quant-backtest} without \crate{quant-lib} in the workspace.
          Confirm that all tests still pass. This verifies the
          non-destructive property.
    \item \textbf{Cycle detection.} Search the source crates for any
          \texttt{use quant\_lib} statement. If you find one, it is a
          bug. Report the file and line.
\end{enumerate}

\section{Summary}

\crate{quant-lib} is the unified library: a
single crate that re-exports nine specialised crates under a unified
namespace. It adds no new logic, depends on every source crate, and
exposes a curated prelude for ergonomic imports. The
\func{GenericBacktest} engine, built on the pluggable trait
catalogue, makes the library composable: swapping a labeler,
validator, or bet sizer is a one-line change via the
\func{BacktestBuilder} fluent API. The 16 examples progress from
single-crate usage to multi-crate pipelines, giving the reader a
guided tour of the ecosystem. The facade is non-destructive: every
source crate continues to work without it, and the 20-test contract
includes a test that verifies this property.